\section{Continuous evidentiary chain of the exceptional publication: from the APP unfolding to Paley--Witt,
\texorpdfstring{\(V^\natural\)}{V natural} and Moonshine}
\label{sec:prueba-corredor-excepcional}

The exceptional publication starts from the incidence reader
\(\mathfrak E_{\rm exc}\) of \eqref{eq:frontal-doble-proyeccion}. It does not take an already given
code, lattice or exceptional group as input: it takes the
boundary, orientation, sheet and memory of an enriched history and
expresses them in the projective gauge fixed by nonadic holonomy. Only
afterwards are the classical theorems of codes, lattices and vertex operator
algebras applied to the produced objects.

\subsection{\texorpdfstring{APP unfolding and graded cone of the index \(54\)}{APP unfolding and graded cone of the index 54}}

Let \(R\) be the involution exchanging the APP sheets and let
\(P_+=(I+R)/2\), \(P_-=(I-R)/2\) be its projectors. The central block
\begin{equation}
 B_c=7I+2R=9P_++5P_-
 \label{eq:exc-bloque-central-app}
\end{equation}
produces four objects of distinct types:
\begin{equation}
 9-5=4,\qquad (B_c)_{i\ne j}=2,\qquad
 \sum M_\Pi-\sum M_\Sigma=51-45=6,\qquad
 9\cdot6=459-405=54.
 \label{eq:exc-despliegue-42654}
\end{equation}
The \(4\) is the spectral jump between sheets, the \(2\) the local coupling,
the \(6\) the difference of the modular compressions and the \(54\) the
integration of that difference over the nine positions. Therefore,
\(4\to2\to6\to54\) designates a hierarchy of morphisms and evaluations, not
a chain of equalities between objects.

With the orientation \(o_{\Sigma\Pi}\), sheet parity and triangular
normalization fixed, the Fock branch has the coupling
\begin{equation}
 \mathcal H_{\Sigma,\Pi}^{(m)}:
 M_{\rm bal}^{C_3,-}\longrightarrow
 \operatorname{End}\!\bigl(\mathcal F_2(\mathcal V_m)\bigr)
 \otimes o_{\Sigma\Pi},
 \label{eq:exc-acoplamiento-fock}
\end{equation}
and its oriented trace satisfies
\begin{equation}
 \operatorname{Tr}_{o,\mathcal F_2}
 \bigl(\mathcal H_{\Sigma,\Pi}^{(m)}(\nu)\Gamma_2(g_m)\bigr)
 =-\frac{3m\,c(\nu)}2.
 \label{eq:exc-traza-fock}
\end{equation}
For \(m=12\) and \(c(\delta)=-3\), \(54\) holds. Independently, the
cyclotomic product
\begin{equation}
 t_3(q)=q^{-1}\prod_{n\ge1}(1+q^n+q^{2n})^{-12}
 =q^{-1}-12+54q-76q^2-243q^3+\cdots
 \label{eq:exc-t3}
\end{equation}
gives \([q]t_3=54\).

\begin{theorem}[Typed cone of the common index]
\label{thm:exc-cono-54}
Once the APP census, the Fock branch, the cyclotomic
product, the lattice neighbor of
\cref{eq:exc-vecino-indice-tres} and the lunar module of
\cref{eq:exc-flm} have been constructed separately, their readers satisfy
\begin{equation}
 D_{\rm APP}
 =\operatorname{Tr}_{o,\mathcal F_2}
  \bigl(\mathcal H_{\Sigma,\Pi}^{(12)}(\delta)\Gamma_2(g_{12})\bigr)
 =[q]t_3=v^2=\frac{196884}{5\cdot729+1}=54.
 \label{eq:exc-cono-54}
\end{equation}
The equality gathers evaluations in a cone toward \(\mathbb Z\); it does not
identify their domains or use \(196884\) to select the corridor.
\end{theorem}

\begin{proof}
The first three equalities are
\eqref{eq:exc-despliegue-42654}, \eqref{eq:exc-traza-fock} and
\eqref{eq:exc-t3}. For the radial vector
\(v=4\rho_1+\rho_2+\cdots+\rho_{12}\) of
\eqref{eq:exc-vector-vecino}, orthogonality of the factors \(A_2\) gives
\(v^2=4^2\cdot2+11\cdot2=54\). Finally,
\(5\cdot729+1=3646\) and \(54\cdot3646=196884\). Each identity is proved in
its own codomain before they are compared.
\end{proof}

\subsection{\texorpdfstring{Paley form and internal quadratic root of \(-I_6\)}{Paley form and internal quadratic root of -I6}}

Let
\[
 \mathcal U_9=(\mathbb Z/9\mathbb Z)^\times=\{1,2,4,8,7,5\},
\]
in the order of the powers of \(2\), and let
\begin{equation}
 \theta:\mathcal U_9\xrightarrow{\sim}\mathbb P^1(\mathbb F_5),
 \qquad
 \begin{array}{c|cccccc}
 u&1&2&4&8&7&5\\ \hline
 \theta(u)&\infty&0&1&2&3&4
 \end{array}.
 \label{eq:exc-calibre-proyectivo}
\end{equation}
This bijection is a coordinate gauge of the reader; it is not used as a
homomorphism. For the quadratic character
\(\chi:\mathbb F_5\to\{-1,0,1\}\), define
\begin{equation}
 C_P=
 \begin{pmatrix}
 0& 1& 1& 1& 1& 1\\
 1& 0& 1&-1&-1& 1\\
 1& 1& 0& 1&-1&-1\\
 1&-1& 1& 0& 1&-1\\
 1&-1&-1& 1& 0& 1\\
 1& 1&-1&-1& 1& 0
 \end{pmatrix},
 \label{eq:exc-matriz-paley}
\end{equation}
where the finite off-diagonal entries are \(\chi(x-y)\) and the
entries incident with \(\infty\) are one.

\begin{lemma}[Quadratic correlation]
\label{lem:exc-correlacion-cuadratica}
If \(a\ne b\) in \(\mathbb F_5\), then
\[
 \sum_{t\in\mathbb F_5}\chi(t-a)\chi(t-b)=-1.
\]
\end{lemma}

\begin{proof}
The substitution \(u=(t-a)/(b-a)\) reduces the sum to
\(\sum_u\chi(u(u-1))\). For \(u=0,1,2,3,4\), the summands are
\(0,0,1,-1,-1\).
\end{proof}

\begin{proposition}[Conference identity and ternary reduction]
\label{prop:exc-paley-witt}
The matrix \(C_P\) is symmetric and satisfies
\[
 C_PC_P^{\mathsf T}=5I_6.
\]
Its reduction \(A_W=C_P\bmod3\) satisfies
\begin{equation}
 \boxed{A_W^2=-I_6\quad\text{in }M_6(\mathbb F_3).}
 \label{eq:exc-aw-cuadrado}
\end{equation}
\end{proposition}

\begin{proof}
Each row of \(C_P\) has squared norm five. Two distinct finite
rows have inner product \(1-1=0\) by the
\cref{lem:exc-correlacion-cuadratica}; the sum of a nontrivial character
also proves orthogonality with the row of \(\infty\). Reducing modulo
three yields
\(A_WA_W^{\mathsf T}=5I_6=2I_6=-I_6\). Since \(A_W\) is symmetric,
\eqref{eq:exc-aw-cuadrado} follows.
\end{proof}

The identity \eqref{eq:exc-aw-cuadrado} constructs the finite complex
structure of the corridor. Indeed,
\[
 \mathbb F_9=\mathbb F_3[\iota]/(\iota^2+1),
 \qquad
 P_{\pm\iota}=\frac12(I\mp\iota A_W)
\]
are complementary rank-three projectors and endow
\(\mathbb F_3^6\) with a rank-three \(\mathbb F_9\)-module structure.
The action is explicit:
\begin{equation}
 (a+b\iota)\cdot v=av+b(vA_W).
 \label{eq:exc-accion-f9}
\end{equation}
Therefore, \(3^6=9^3\) is not used as a cardinal substitute for an action.

\subsection{Self-dual ternary code and Witt system}

Define
\begin{equation}
 c:\mathbb F_3^6\longrightarrow\mathbb F_3^{12},
 \qquad c(q)=(q,qA_W),
 \qquad C_W:=\operatorname{im}c.
 \label{eq:exc-codigo-grafico}
\end{equation}

\begin{theorem}[Ternary code produced by incidence]
\label{thm:exc-codigo-ternario}
The map \(c\) is injective and
\begin{equation}
 C_W=C_W^\perp,
 \qquad
 W_{C_W}(y)=1+264y^6+440y^9+24y^{12}.
 \label{eq:exc-enumerador-cw}
\end{equation}
In particular, \(C_W\) has parameters \([12,6,6]_3\).
\end{theorem}

\begin{proof}
The first projection of \(c(q)\) is \(q\), so \(c\) is injective.
With \(G_W=(I_6\mid A_W)\),
\[
 G_WG_W^{\mathsf T}=I_6+A_WA_W^{\mathsf T}=0
\]
in \(\mathbb F_3\). Thus \(C_W\subseteq C_W^\perp\), and equality of
dimensions gives self-duality. The enumerator is the exact finite sum
\[
 \sum_{q\in\mathbb F_3^6}
 y^{\operatorname{wt}(q)+\operatorname{wt}(qA_W)}.
\]
Evaluating the \(3^6=729\) rows of the explicit matrix
\eqref{eq:exc-matriz-paley} gives, respectively, the counts
\(1,264,440,24\) at weights \(0,6,9,12\), and zero at the others. The sum
of the four counts is \(729\), so no word remains outside
the finite certificate.
\end{proof}

The \(264\) weight-six vectors occur in pairs \(\{x,-x\}\); their
\(132\) supports form the Steiner system
\begin{equation}
 W_{12}=S(5,6,12).
 \label{eq:exc-w12}
\end{equation}
This is the recognition, after constructing \(C_W\), of the extended ternary
Golay code and its Witt design
\citep{MacWilliamsSloane1977,ConwaySloane1999}. The design is not introduced
to select the matrix \(A_W\).

The binary prolongation remains separate. If
\(\mathcal G_{24}\subset\mathbb F_2^{24}\) is the extended binary
Golay code, \(D\) a dodecad and
\(\ell:\{1,\ldots,12\}_{\rm HMT}\to D\) a design isomorphism, the
weight-eight supports form \(W_{24}=S(5,8,24)\). For a fixed tetrad,
\[
 \lambda_4(W_{12})=4,
 \qquad
 \lambda_4(W_{24})=5,
\]
so that four octads lift the four incident hexads and a fifth
is transversal. The datum \((\mathcal G_{24},D,\ell)\) belongs to the binary
codomain; there is no canonical arrow \(W_{24}\to\Lambda_{24}\).

\subsection{Ternary gluing and Leech neighbor}

Let \(Q=A_2^{12}\), with \(\det Q=3^{12}\). The discriminant of each
factor \(A_2\) is \(\mathbb F_3\), and the gluing defined by
\eqref{eq:exc-codigo-grafico} is
\begin{equation}
 N:=\{x\in Q^*:x\bmod Q\in C_W\}.
 \label{eq:exc-niemeier-pegado}
\end{equation}

\begin{theorem}[Niemeier lattice produced by the code]
\label{thm:exc-niemeier}
The lattice \(N\) is positive definite, even and unimodular of rank
twenty-four, and its root system is \(A_2^{12}\). Consequently,
\(N\cong N(A_2^{12})\).
\end{theorem}

\begin{proof}
Self-duality of \(C_W\) makes the gluing integral. All its weights are
multiples of three, which makes the extension even. Moreover,
\([N:Q]=|C_W|=3^6\), and therefore
\[
 \det N=\frac{\det Q}{[N:Q]^2}
 =\frac{3^{12}}{3^{12}}=1.
\]
A nonzero discriminant class of \(A_2\) has minimum norm \(2/3\).
Since every nonzero word of \(C_W\) has weight at least six, every gluing
vector outside \(Q\) has norm at least four. The norm-two
vectors are exactly the roots of the twelve factors \(A_2\). The
classification of even unimodular lattices of rank twenty-four
identifies the indicated type \citep{ConwaySloane1999}.
\end{proof}

Let \(\rho_i\) be the highest root of the factor \(A_2^{(i)}\), normalized by
\(\rho_i^2=2\), and define
\begin{equation}
 v=4\rho_1+\rho_2+\cdots+\rho_{12},
 \qquad v^2=54.
 \label{eq:exc-vector-vecino}
\end{equation}
The ordered frame of the exceptional reader fixes which of the twelve fibers occupies
the first position; no Archimedean value intervenes in this selection.
Set
\begin{equation}
 \chi_v(x)=\langle x,v\rangle\bmod3,
 \quad N_0=\ker\chi_v,
 \quad N_v=N_0+\mathbb Z\frac v3.
 \label{eq:exc-vecino-indice-tres}
\end{equation}

\begin{lemma}[Structure of the neighbor]
\label{lem:exc-estructura-vecino}
The character \(\chi_v\) is surjective, \([N:N_0]=3\),
\(v/3\in N_0^*\), and \(N_v\) is even and unimodular.
\end{lemma}

\begin{proof}
For every root of a factor \(A_2\), the product with the highest root is
\(\pm1\) or \(\pm2\); the coefficients of \(v\) are congruent to one
modulo three. Therefore, no root belongs to \(N_0\) and \(\chi_v\) is
nontrivial, hence surjective. The definition of the kernel gives index three and
\(\det N_0=9\). For \(x\in N_0\),
\(\langle x,v/3\rangle\in\mathbb Z\), while
\((v/3)^2=6\). Thus the extension by \(v/3\) is integral and even. Since
\([N_v:N_0]=3\),
\(\det N_v=9/3^2=1\).
\end{proof}

\begin{lemma}[Minimum of the new classes]
\label{lem:exc-minimo-clases}
Each coordinate of a local class appearing in
\(N_0\pm v/3\) has minimum norm \(2/9\). Consequently,
\[
 \|x\pm v/3\|^2\ge12\cdot\frac29=\frac83
 \qquad(x\in N_0).
\]
\end{lemma}

\begin{proof}
Writing a coordinate of \(A_2^*\) in the simple-root basis, the
six relevant oriented classes are the shifts of
\(\pm\rho/3\) by the three discriminant classes. Completing squares in the
Gram form
\(\bigl(\begin{smallmatrix}2&-1\\-1&2\end{smallmatrix}\bigr)\)
gives minimum \(2/9\) in all six classes. The orthogonal sum of the twelve
coordinates gives the inequality.
\end{proof}

\begin{theorem}[Construction of the Leech lattice]
\label{thm:exc-leech}
The neighbor \(N_v\) of \eqref{eq:exc-vecino-indice-tres} is isometric to the
Leech lattice:
\begin{equation}
 \boxed{N_v\cong\Lambda_{24}.}
 \label{eq:exc-leech}
\end{equation}
\end{theorem}

\begin{proof}
The roots of \(N\) do not belong to \(N_0\), and the two new classes
contain no norm-two vectors by the
\cref{lem:exc-minimo-clases}. The \cref{lem:exc-estructura-vecino} proves that
\(N_v\) is positive definite, even, unimodular and of rank twenty-four. The
uniqueness of the lattice with these properties and no roots identifies it with
\(\Lambda_{24}\) \citep{ConwaySloane1999}.
\end{proof}

\subsection{Order-three isometry, orbifold and lunar module}

The Coxeter rotation of each factor \(A_2\) induces an order-three isometry
\(g_c\). Stability of the code \(C_W\) proves that
\(g_cN=N\). For the oriented frame fixing \(v\), the shifts
\begin{equation}
 y_*:=\frac{g_cv-v}{3},
 \qquad
 z_*:=\frac{g_c^{-1}v-v}{3}
 \label{eq:exc-desplazamientos-vecino}
\end{equation}
belong to \(N_0\): their discriminant words are a pair
\(c_*,-c_*\in C_W\), and
\(\langle y_*,v\rangle=\langle z_*,v\rangle=-27\). Therefore,
\begin{equation}
 g_cN_v=N_v.
 \label{eq:exc-preservacion-vecino}
\end{equation}
The action has no fixed points in \(N_v\otimes\mathbb R\), since each
factor \(A_2\) has the two primitive eigenvalues of order three.

Let \(\widehat g_c\) be the standard lift to the lattice VOA
\(V_{\Lambda_{24}}\). Its two nontrivial eigenspaces have dimension
twelve. The vacuum energy of the twisted module is therefore
\begin{equation}
 \rho_{g_c}
 =\frac1{4\cdot3^2}
 \sum_{k=1}^{2}k(3-k)\dim\mathfrak h_{(k)}
 =\frac43.
 \label{eq:exc-peso-torcido}
\end{equation}
The congruence \(3^2\rho_{g_c}=12\equiv0\pmod3\) proves that the lift
is of type zero.

\begin{theorem}[Triadic orbifold]
\label{thm:exc-orbifold}
The cyclic orbifold of the preceding lift satisfies
\begin{equation}
 \operatorname{Orb}_{\mathbb Z/3}
 (V_{\Lambda_{24}},\widehat g_c)\cong V^\natural.
 \label{eq:exc-orbifold}
\end{equation}
The dual automorphism belongs to the class \(3B\).
\end{theorem}

\begin{proof}
The three hypotheses that cannot be replaced by an equality of series
have been verified: fixed-point-free order-three isometry, preservation of the
neighbor and type zero. The cyclic orbifold theorems then identify
the holomorphic extension with the lunar module and classify the dual automorphism
\citep{ChenLamShimakura2016,AbeLamYamada2017,Carnahan2017}.
\end{proof}

The order-two construction is a parallel realization. For the lift
\(\theta\) of negation on \(\Lambda_{24}\),
\begin{equation}
 V^\natural=(V_{\Lambda_{24}})^+
 \oplus(V_{\Lambda_{24}}^T)^+.
 \label{eq:exc-flm}
\end{equation}
The order-two and order-three routes are identified by the VOA theorems; they are not
causally serialized one inside the other \citep{FLM1988}.

\begin{theorem}[Publications of the lunar module]
\label{thm:exc-moonshine}
Once \(V^\natural\) has been constructed, one obtains
\begin{align}
 \operatorname{Aut}(V^\natural)&\cong\mathbb M,
 \label{eq:exc-aut-monster}\\
 \operatorname{ch}_{V^\natural}(\tau)
 &=J(\tau)=j(\tau)-744
 =q^{-1}+196884q+21493760q^2+\cdots,
 \label{eq:exc-caracter-j}\\
 T_g(\tau)&=\operatorname{Tr}_{V^\natural}
 \bigl(gq^{L_0-1}\bigr),\qquad g\in\mathbb M.
 \label{eq:exc-mckay-thompson}
\end{align}
The series \(T_g\) are the genus-zero modular functions of the Moonshine
theorem and satisfy the replicability identities.
\end{theorem}

\begin{proof}
The first identification and the graded construction come from the
FLM realization. The coordination between the graded action, the McKay--Thompson
series and their modular properties is the Moonshine theorem of
Conway--Norton and Borcherds
\citep{FLM1988,ConwayNorton1979,Borcherds1992}. In particular,
\((V^\natural)_0=\mathbb C\mathbf1\) and
\((V^\natural)_1=0\); the coefficient \(196884=1+196883\) is the first
nontrivial decomposition into dimensions of Monster representations.
\end{proof}

\begin{falsifier}[Exceptional corridor]
\label{fals:exc-corredor}
The publication is refuted if any of the effective steps fails:
\eqref{eq:exc-despliegue-42654}, the trace
\eqref{eq:exc-traza-fock}, the coefficient \([q]t_3=54\),
\(C_PC_P^{\mathsf T}=5I_6\), \(A_W^2=-I_6\), self-duality or the
enumerator of \(C_W\), evenness or unimodularity of the gluing, absence
of roots of the neighbor, its preservation by \(g_c\), the twisted weight
\(4/3\) or type zero. A narrative is also refuted if it uses
\(W_{24}\) to construct \(\Lambda_{24}\), confuses the two sets
of \(243\) on account of their cardinality, or makes the Monster act directly
on digits and routes. The action of \(\mathbb M\) resides in \(V^\natural\).
\end{falsifier}
